\section{An integrable boundary for the complete-intersection argument}
\label{sec:ci-cross-minor-boundary}

\paragraph{Status and scope.}
The full derivation and exact supplement have passed the root check and
an independent whole-proof check by /root/sixfolds. This is a
counterexample to an overstrong scalar inference. It supplies neither a
module deformation nor a weakly semiregular object and does not reduce
the main supplier gate.

\begin{proposition}[HD-CI-B1: scalar cross minors do not force a zero Jacobian]
\label{HD-CI-B1}
There is a length-sixteen complete intersection
$A=\mathbf C[[x,y,z,w]]/(f_1,f_2,f_3,f_4)$ of embedding dimension four
with complementary constant bivectors
$\pi_+=\partial_x\wedge\partial_y$ and
$\pi_-=\partial_z\wedge\partial_w$ such that all scalar minors
\[
 P^+_{ij}P^-_{kl}-P^-_{ij}P^+_{kl}=0\quad\hbox{in }A,
 \qquad P^{\pm}_{ij}=\{f_i,f_j\}_{\pi_\pm},
\]
while $\det\operatorname{Jac}(f)$ is a nonzero socle element.
No regular generating sequence for this ideal has the commuting-triple
property of SF-E836 with these same two bivectors.
\end{proposition}

\begin{proof}
Let $J$ be the rational matrix with successive columns
\[
 u=(1,0,-3/5,-8/5)^t,\quad v=(0,1,1,1)^t,\quad
 u'=(1,0,-9/5,-16/5)^t,\quad v'=(0,1,2,3)^t.
\]
In the coordinate order $(12,13,14,23,24,34)$, the two column wedges are
\[
 p=u\wedge v=(1,1,1,3/5,8/5,1),\qquad
 q=u'\wedge v'=(1,2,3,9/5,16/5,1).
\]
Both are decomposable. The ratios $q_{ij}/p_{ij}$ agree within each
complementary pair, with values respectively $1,2,3$. Hence every
complementary scalar coefficient minor vanishes. Nevertheless
\[
 p\wedge q=2(1-16/5+9/5)=-4/5=\det J\ne0.
\]
Put $\ell=J(x,y,z,w)^t$ and $f_i=\ell_i^2$. Invertibility gives
\[
 A\simeq\mathbf C[[\ell_1,\ell_2,\ell_3,\ell_4]]/
                 (\ell_1^2,\ell_2^2,\ell_3^2,\ell_4^2).
\]
The sixteen squarefree monomials are a basis, so the asserted length
and embedding dimension follow. These are actual power series, whose
actual gradient matrix is
\[
 \operatorname{Jac}(f)=2\operatorname{diag}(\ell_1,\ell_2,\ell_3,\ell_4)J.
\]
It follows that
$P^+_{ij}=4\ell_i\ell_jp_{ij}$ and
$P^-_{ij}=4\ell_i\ell_jq_{ij}$.
If the pairs $(ij)$ and $(kl)$ share an index, both terms in the scalar
minor contain a square $\ell_a^2$ and vanish in $A$.
If they are disjoint, the scalar coefficient minor vanishes by the
explicit complementary-pair calculation. This proves all the stated
minor identities, including coincident and ordered pairs.

The actual Jacobian determinant is
\[
 \det\operatorname{Jac}(f)=-\frac{64}{5}\ell_1\ell_2\ell_3\ell_4.
\]
The top squarefree monomial is nonzero and is annihilated by the maximal
ideal, so it is a nonzero socle element. The coefficient of the mixed
bracket wedge is the same nonzero element.

For completeness, replacing a minimal regular generating sequence by
another multiplies its column of generators by an invertible matrix
$U$ modulo the ideal. Modulo that ideal the derivative terms involving
$U$ disappear, so both bracket arrays transform by $\bigwedge^2U$.
The vanishing of all scalar minors is preserved, while their mixed
wedge is multiplied by $\det U$, a unit. If a new generator triple
commuted under $\pi_+$, each surviving mixed-wedge term would be one
of those vanishing scalar minors. This would contradict its invariant
nonvanishing. Thus no such generator presentation exists.
\end{proof}

\paragraph{Logical boundary and exact checks.}
The two companion matrix orders of SF-E836 imply the displayed scalar
minor identities, but HD-CI-B1 does not solve the companion equations.
It does not construct a flat deformation of any coherent hull-modification
module. The root and separate peer programs check all thirty-six ordered
minor equations, all actual-gradient entries, and the nonzero socle;
these checks supplement the complete two-case proof above.
In the common Grothendieck framework this is a local Artin scheme and
$K$-theoretic deformation boundary for a proposed condition-A supplier
(cote16, scan050; SIGN\_NEXT SN-C03/C04), not a condition-I production
of cycles (scan049, cote44, Theorem4). No change to the main detector
gate or to the conditional status of Markman, Severn or032 follows.
